\section*{Реферат}
\addcontentsline{toc}{section}{Реферат}

Пояснительная записка содержит 49 страниц текста, 5
рисунков, 6 таблиц, 2 приложения, список из 23 источников.

Ключевые слова: непрерывное профилирование, CI/CD, регрессия
производительности, Go, pprof, benchstat, бенчмаркинг, статистическая
значимость, критерий Уэлча, автоматизация тестирования.

Работа посвящена разработке и экспериментальной апробации прототипа
системы, интегрирующей профилирование производительности Go-приложений в
CI/CD-пайплайн для раннего обнаружения регрессий производительности —
до попадания кода в основную ветку репозитория. Пояснительная записка
состоит из пяти основных разделов: аналитического (обзор инструментов
профилирования Go и подходов к их интеграции в CI/CD), теоретического
(формализация критерия обнаружения регрессии и модели процесса
профилирования), инженерно-технологического (проектирование архитектуры
прототипа), практического (программная реализация прототипа и CI-пайплайна)
и раздела экспериментальной апробации (проверка корректности обнаружения
регрессий и оценка накладных расходов профилирования на четырёх
управляемых сценариях). Работа завершается заключением с оценкой степени
достижения поставленной цели и направлений дальнейших исследований.
